% GENERATED, edit the French unit or the translation table
\setcounter{chapter}{11}\mxunit{terminale-c13}{Parametric representations and Cartesian equations}\label{ch:12}

{\small\sffamily\color{mxmint}S65, S67, S70\par}

\begin{keypoint}[Objectives]Describe a line and a plane; determine intersections; choose a method for a projection.\end{keypoint}

\section{Course notes}
Throughout this chapter, coordinates are given in an orthonormal frame of space. This assumption is necessary to identify a normal vector from the coefficients of an equation and to calculate distances.
\subsection{Parametrised lines}
The line through $A(x_A,y_A,z_A)$ with non-zero direction vector $(a,b,c)$ has the representation
\[d:\quad\begin{cases}x=x_A+at\\y=y_A+bt\\z=z_A+ct\end{cases}\quad(t\in\R).\]
A point belongs to the line if there exists \emph{one and the same} real number $t$ satisfying the three equalities. A line has many representations: change the starting point or multiply the direction vector by a non-zero real number.
\paragraph{Example.} $(x,y,z)=(1,2,0)+t(2,-1,3)$ describes the same line as $(3,1,3)+s(4,-2,6)$; the link is $t=1+2s$.
\subsection{Cartesian planes}
\begin{keypoint}[Equation of a plane]
If $(a,b,c)\ne(0,0,0)$, the set $ax+by+cz+d=0$ is a plane with normal vector $(a,b,c)$. The plane with normal $\vec n$ through $A$ is constructed by writing $\vect{AM}\cdot\vec n=0$.
\end{keypoint}
\paragraph{Proof.} The condition $\vect{AM}\cdot\vec n=0$ is written $a(x-x_A)+b(y-y_A)+c(z-z_A)=0$. Expanding, $d=-ax_A-by_A-cz_A$. Conversely, every point satisfying this equation satisfies the orthogonality condition, hence belongs to the plane.
\paragraph{Example.} The plane with normal $(2,-1,3)$ through $A(1,2,0)$ is $2(x-1)-(y-2)+3z=0$, that is $2x-y+3z=0$.
\subsection{Intersections}
For a line and a plane, replace $x,y,z$ in the equation of the plane by their parametric expressions. A solution in $t$ gives a single intersection point; an identity gives a line contained in the plane; a contradiction gives a parallel line that is disjoint from the plane.
For two planes, non-collinear normals give an intersection along a line; solve the system of two equations by choosing one variable as a parameter. If the normals are collinear, the planes are parallel: distinguish coincident from disjoint.
For two lines, use independent parameters $s,t$ and solve the three equations. With no solution and non-collinear direction vectors, the lines are not coplanar.
\subsection{Projection onto a plane}
If $\mathcal P:ax+by+cz+d=0$, look for the orthogonal projection of $M$ onto the normal line through $M$: $H=M+t(a,b,c)$. Substitute into the equation of the plane to determine $t$, then calculate $MH$.
We deduce, without making it a formula to be learned in isolation,
\[d(M,\mathcal P)=\frac{|ax_M+by_M+cz_M+d|}{\sqrt{a^2+b^2+c^2}}.\]
\paragraph{Essential check.} After calculating a projection, check simultaneously that $H\in\mathcal P$ and that $\vect{MH}$ is collinear with the normal.


\section{Research activity}\textbf{Three lines facing a plane --- 30 min.}\par In an orthonormal frame of space, consider the plane $\mathcal P$ with equation $x+y+z=3$. In each question, the parameter $t$ ranges over $\R$.
\begin{enumerate}\item The line $d_1$ has the representation $(x;y;z)=(t;t;t)$. Determine $d_1\cap\mathcal P$.
\item The line $d_2$ has the representation $(x;y;z)=(3+t;-t;0)$. Determine $d_2\cap\mathcal P$.
\item The line $d_3$ has the representation $(x;y;z)=(t;-t;0)$. Determine $d_3\cap\mathcal P$.
\item For each line, state whether the equation obtained after substitution into that of the plane has a unique solution, no solution, or every real number as a solution. Link these cases to the geometric positions found.\end{enumerate}

\section{Progressive exercises}

\mxexerciseitem{1}{Practice}In a frame of space, consider $A(1;-1;2)$ and $B(3;0;1)$. Give a parametric representation of the line $(AB)$.

\mxexerciseitem{2}{Practice}In an orthonormal frame, determine a Cartesian equation of the plane through $A(1;2;3)$ with normal vector $\vec n=(1;0;-2)$.

\mxexerciseitem{3}{Practice}In a frame of space, the line $d$ is given by $(x;y;z)=(t;1+t;2-t)$, $t\in\R$, and the plane $\mathcal P$ by $x+2y+z=8$. Determine their intersection.

\mxexerciseitem{4}{Practice}In a frame of space, determine the intersection of the planes with equations $x+y+z=1$ and $x-y=0$. Give a parametric representation if the intersection is a line.

\mxexerciseitem{5}{Practice}In a frame of space, study the relative position of the planes with equations $2x-y+z=3$ and $4x-2y+2z=5$. Justify.

\section{Textbook or past-paper work}\textbf{Baccalauréat, Metropolitan France, 17 June 2025: exercise 3, statements 1 to 4 (p.~5--6).} At home after S70. Justify each answer; for a false statement, explain precisely what refutes it. The paper is available here: \href{https://www.education.gouv.fr/sites/default/files/2025-06/baccalaur-at-g-n-ral-2025-math-matiques-440856.pdf}{official past paper}. Use the data and instructions of this past paper.

\section{Homework}\textbf{Homework 12 --- Projection and intersection --- 55 min.}\par In an orthonormal frame of space, consider the point $M(3;0;0)$ and the planes $\mathcal P:x+2y+2z=9$ and $\mathcal Q:x-y=0$.
\begin{enumerate}\item Give the coordinates of a non-zero normal vector to $\mathcal P$.
\item Determine the coordinates of the orthogonal projection $H$ of $M$ onto $\mathcal P$, then calculate the distance from $M$ to $\mathcal P$.
\item Show that $\mathcal P\cap\mathcal Q$ is a line and give a parametric representation of it.
\item Does the point $H$ belong to this line? Justify.
\item Determine the coordinates of the point $K$ of this line whose third coordinate (the ``cote'') is zero. Calculate $MK$, then compare $MK$ and $MH$.\end{enumerate}

\ifmxteacher\mxnewpage\section{Teaching notes and solutions}\begin{teachernote}[Teacher notes]Use two different parameters to compare two lines. Check all three coordinates, even if two suffice to calculate candidates. A Cartesian equation on its own generally describes a plane in space; it does not describe a line. Reuse the projection from Chapter 11 to give meaning to the solving of systems.\end{teachernote}

\subsection{Solutions to the activity}\begin{enumerate}
\item A point of $d_1$ has coordinates $(t,t,t)$. It belongs to the plane if and only if $3t=3$, that is $t=1$. The intersection is therefore the single point $(1,1,1)$.
\item For a point of $d_2$, $x+y+z=3+t-t+0=3$ for every $t$. All points of the line are in the plane: $d_2\subset\mathcal P$.
\item For a point of $d_3$, $x+y+z=t-t+0=0$. The equation $0=3$ has no solution: the line is disjoint from the plane. Its direction vector $(1,-1,0)$ is orthogonal to the normal $(1,1,1)$; it is parallel to the plane.
\item After substitution, we obtain respectively an equation with a single solution, an identity true for every parameter value, and a contradiction. These three cases correspond to a single intersection point, containment, and no intersection.
\end{enumerate}

\paragraph{E1.}$\vect{AB}=(2,1,-1)\ne\vec0$. The line therefore has the representation $(x,y,z)=(1,-1,2)+t(2,1,-1)$, $t\in\R$. The values $t=0$ and $t=1$ give back $A$ and $B$.

\paragraph{E2.}For $A(1,2,3)$, the condition $\vect{AM}\cdot(1,0,-2)=0$ is $(x-1)-2(z-3)=0$. After expanding, the plane has equation $x-2z+5=0$. The point $A$ satisfies it.

\paragraph{E3.}After substitution, $t+2(1+t)+(2-t)=8$, that is $2t+4=8$, so $t=2$. The unique intersection is $(2,3,0)$; we check $2+2\times3+0=8$.

\paragraph{E4.}The second equation gives $x=y$. Let $x=y=t$; the first becomes $2t+z=1$. The set of solutions is $(x,y,z)=(t,t,1-2t)=(0,0,1)+t(1,1,-2)$, a line.

\paragraph{E5.}The normals are collinear, the second being twice the first. But doubling $2x-y+z=3$ gives $4x-2y+2z=6$, which is incompatible with the right-hand side $5$. The planes are therefore strictly parallel, not intersecting.

\mxneed{30}

\subsection{Solutions to Homework 12}\begin{enumerate}
\item The coefficients of $x,y,z$ give the non-zero normal $\vec n=(1,2,2)$ in the orthonormal frame.
\item Look for $H$ on the normal line through $M$: $H=(3+t,2t,2t)$. Membership of $\mathcal P$ gives $3+t+4t+4t=9$, so $t=\frac{2}{3}$. Thus $H=(\frac{11}{3},\frac{4}{3},\frac{4}{3})$. We check $\frac{11}{3}+\frac{8}{3}+\frac{8}{3}=9$ and $\vect{MH}=(\frac{2}{3})\vec{n}$. So $H$ is indeed the orthogonal projection, and $MH=(\frac{2}{3})\sqrt{1+4+4}=2$.
\item The normals $(1,2,2)$ and $(1,-1,0)$ are not collinear: the planes meet along a line. Since $x=y$, let $x=y=s$. The first equation gives $3s+2z=9$, that is $z=\frac{9-3s}{2}$. The intersection is the line $(x,y,z)=(s,s,\frac{9-3s}{2})$, $s\in\R$.
\item The point $H$ has $x_H=\frac{11}{3}\ne\frac{4}{3}=y_H$, so it belongs neither to $\mathcal Q$ nor to the line of intersection.
\item The cote is the $z$-coordinate. Setting it to zero requires $9-3s=0$, so $s=3$ and $K=(3,3,0)$. Then $MK=3>2=MH$. This result is consistent with the property of the orthogonal projection: among all points of the plane, $H$ is the closest to $M$.
\end{enumerate}

\subsection{Solutions to the textbook or past-paper work}\begin{enumerate}
\item True: the line passes through $B$ and its direction vector $(-2;-1;3)$ equals $-\frac12\vect{AB}$.
\item False: $\vec n\cdot\vect{OB}=15-4-1=10\ne0$.
\item False: for $k=-4$ and $s=\frac52$, the two representations give $(11;12;-14)$. The lines intersect, so they are coplanar.
\item True: $d(C,\mathcal P)=\frac{|2-(-1)+2+1|}{\sqrt{1+1+1}}=2\sqrt3$.
\end{enumerate}

\fi

